\subsection{Published counts and the exact scopes of our catalogs}

\begin{proposition}[Proved: catalog completeness in the scopes actually used]
\label{CP2:A-catalog-scope-comparison-through-ten}
Our ordinary monoid catalogs are complete in every order from one
through seven. Our source catalogs in orders eight, nine and ten
are complete for the condition $H_2(BM;\mathbb Q)\ne0$; they are
not catalogs of all monoids in those orders. The exact counts are
\[
\begin{array}{c|r|r}
\text{order}&\text{complete ordinary catalog}&H_2(\mathbb Q)\ne0\\\hline
1&1&0\\
2&2&0\\
3&7&0\\
4&35&0\\
5&228&1\\
6&2237&5\\
7&31559&43
\end{array}
\]
\[
\begin{array}{c|r|r}
\text{order}&\text{published ordinary count}&\text{positive source catalog}\\\hline
8&1668997&510\\
9&3685886630&8046\\
10&\text{not listed}&157145
\end{array}
\]
Here ``not listed'' refers only to the opened ordinary-count source,
not to a claim that the number is unknown. The complete target
catalogs restricted to at least two minimal rows and two minimal
columns have $1205$ entries in order eight and $21920$ in order nine.
Those target catalogs include their rational-$H_2$-zero entries.
\end{proposition}
\begin{proof}
Theorem~\ref{CP2:A-monoid-census} gives the independent full table
enumeration, unrestricted-permutation checker and exact equality
with the supplied RESIDUE census in orders one through five.
Its counts agree in each order with the official OEIS sequence
\cite{OEIS_A058129}; the agreement supplements the exhaustion
proof and does not replace it. The positive subsets through order
seven are compared as complete canonical sets with those full
catalogs. Proposition~\ref{CP2:A-positive-corners-eight},
Proposition~\ref{CP2:A-positive-corners-nine} and
Theorem~\ref{CP2:A-positive-corners-ten} prove the complete
positive-source scopes in the larger orders, including the two
nontrivial-minimal-group classes in order nine and fourteen in
order ten. Their proof reductions are what justify omitting other
monoids from a source search.

The order-eight two-sided target count is
$1181+12+12=1205$. In order nine it is
$21534+191+191+1+1+2=21920$, by the complete minimal-ideal
families in the order-nine proposition. Positivity is not imposed
on either target list. The stronger jointly oriented
eighteen-arrow reduction in
Corollary~\ref{CP2:A-two-object-eighteen-oriented-twenty-two}
subsequently removes the profiles requiring those larger targets;
the earlier catalogs and their exact scopes remain valid.

For comparison, the opened Distler--Kelsey university repository
abstract reports the following counts up to \emph{isomorphism and
anti-isomorphism}, also listed in OEIS A058133:
\[
\begin{array}{c|r}
\text{order}&\text{isomorphism--opposition classes}\\\hline
8&858977\\
9&1844075697\\
10&52991253973742.
\end{array}
\]
We use that abstract only for its published counts; we do not claim
to have read the article's enumeration proof
\cite{CP2-DistlerKelsey2009,OEIS_A058133}.
If $a$ is the ordinary count, $b$ the opposition-orbit count and
$s$ the number of self-opposite ordinary classes, opposition is
an involution and $a=2b-s$. The opened A058132 values
$s_8=48957$ and $s_9=2264764$ give exactly the two ordinary
counts in the table above \cite{OEIS_A058132}. Neither the
ordinary nor the opposition total is equated with a positive-source
subset. The opened ordinary and self-opposite sequences do not
list an order-ten value, so we infer none from them.
\end{proof}

\begin{remark*}[Inherited covers and the completeness obligation]
The inherited order-seven list of $32823$ oriented tables is a
redundant cover, not an ordinary-isomorphism census of size
$31559$. The inherited order-eight positive-source cover has
$516$ entries, including six explicitly marked rational-$H_2$-zero
entries. Our exact positive count is $510$. We neither identify
those list lengths with the published all-monoid totals nor infer
completeness from their agreement. Each finite gate uses the full
catalog or the proved restricted catalog appropriate to its role;
missing historical output is never a premise of this comparison.
\end{remark*}
